\section{The explicit main term}\label{sec:statements}
We first give the counting scale and the main theorem, then construct its
arithmetic coefficient. Definitions involving~$x$ are used only for
sufficiently large~$x$.

\subsection{The scale}
For $j\ge1$, put
\begin{equation}\label{eq:aj}
 a_j=(j+1)\log(j+1)-j\log j-1
     =\int_j^{j+1}\log t\dd t.
\end{equation}
These numbers are positive. The increasing function $\sum_{j\ge1}a_jz^j$ goes from zero to infinity on $0<z<1$, so there is a unique $\rho\in(0,1)$ such that
\[
 \sum_{j\ge1}a_j\rho^j=1.
\]
These coefficients and this root occur in Maier and Pomerance's
asymptotic estimate \cite[Section~1]{MaierPomerance1988}.
Define
\begin{equation}\label{eq:renewal}
 \lambda=\log(1/\rho),\qquad
 \gamma=\left(\sum_{j\ge1}ja_j\rho^j\right)^{-1},\qquad
 g_0=1,\quad g_j=\sum_{d=1}^j a_dg_{j-d}.
\end{equation}
Our $\gamma$ is the constant denoted by $\lambda$ in Ford's recurrence estimates \cite[Section~3]{FordTotients2013}. Set
\[
 C_0=\frac1{2\lambda},\qquad
 D_0=\frac{1+\log(\lambda/\gamma)}{\lambda}-\frac12;
\]
these are the constants in the classical estimate
\eqref{eq:classical-order}. For the scale used in our theorem, put
\begin{equation}\label{eq:scales}
 B=\log_2x,\qquad
 m=\left\lfloor\frac{\log B-\log_2B}{\lambda}\right\rfloor,
 \qquad
 \theta=\frac{\log B-\log_2B}{\lambda}-m,
\end{equation}
so that $0\le\theta<1$, and write
\begin{equation}\label{eq:Gj}
 G_j=\frac{B^j}{j!\prod_{i=1}^j g_i}\quad(0\le j\le m),\qquad G_0=1.
\end{equation}
The factor $xG_m/\log x$ is Ford's counting scale. The following theorem
resolves its bounded multiplicative uncertainty by an explicitly constructed
function of the phase~$\theta$.

\begin{theorem}\label{thm:main}
For sufficiently large positive integers $H$, there are nonnegative
functions $A_H(1;\cdot):[0,1)\to\mathbb R$, given by the finite
arithmetic formula \eqref{eq:AH} below,
which converge uniformly as $H\to\infty$ to a function $A(1;\cdot)$
satisfying
\[
 0<\inf_{0\le s<1}A(1;s)
   \le\sup_{0\le s<1}A(1;s)<\infty.
\]
With the parameters \eqref{eq:renewal}--\eqref{eq:Gj},
\[
 V(x)\sim\frac{x}{\log x}\,G_m A(1;\theta).
\]
Moreover, for every fixed real $c>0$,
\[
 \frac{V(cx)}{V(x)}\longrightarrow c.
\]
\end{theorem}

The first argument of $A$ records a weight; the main theorem uses the
constant weight~$1$. Its arithmetic construction below does not use~$V$.
The fixed-scale conclusion has a separate mechanism: it compares the same
representation count at two endpoints. No regularity of
$s\mapsto A(1;s)$ is asserted or needed.

\subsection{The arithmetic coefficient}\label{subsec:coefficient}
The tail index runs in the opposite direction from the ordered prime
index. In the proof we write a preimage as
$p_0p_1\cdots p_La$, with the primes decreasing, and use two cuts
\[
 R=m-H,\qquad L=m-P,\qquad P=\lfloor\log\log H\rfloor.
\]
The long prefix ends at $p_R$, whereas the exact arithmetic tail
contains $p_{R+1},\ldots,p_L$ and $a$. Thus the notation below is
$Q_h=p_{m-h}$: its indices run from $P$ to $H-1$, and $Q_P$ is the
smallest retained prime. For fixed $H$, these tail primes are bounded
independently of $x$. We first let $x$ tend to infinity with $H$ fixed,
and only afterwards let $H$ tend to infinity. The following definition
uses the phase $s$ directly, without reference to $x$.

Fix $s\in[0,1)$ and put $\alpha_s=\lambda\e^{\lambda s}$. Let $H$ be a sufficiently large positive integer and $P=\lfloor\log\log H\rfloor$. A \emph{tail witness} is a choice
\[
 \eta=((Q_h)_{P\le h<H},a)
\]
of primes $Q_h$ and a positive integer $a$ with the following properties. Writing $v_h=\log_2Q_h$, we require
\begin{equation}\label{eq:tail-bands}
 \begin{aligned}
 0.9\alpha_s h\rho^{-h}&\le v_h\le1.1\alpha_s h\rho^{-h}
          &&(P\le h<H),\\
 \sum_{l=P}^{h-1}a_{h-l}v_l
   &\le(1+10^{-4}\e^{-h/40})v_h
          &&(P\le h<H),\\
 P^+(a)&\le Q_P,\qquad
 \log a\le\exp(2\alpha_sP\rho^{-P}).
 \end{aligned}
\end{equation}
Here $P^+(a)$ is the largest prime factor of~$a$, with $P^+(1)=1$; an empty sum is zero. Associate to~$\eta$ the integer
\[
 w(\eta)=a\prod_{h=P}^{H-1}Q_h.
\]
For $d\in\mathcal V$, let $\mathcal W_{H,s}(d)$ be the set of tail witnesses for which $\varphi(w(\eta))=d$. For $h\ge H$, define
\begin{equation}\label{eq:tail-D}
 D(h,\eta)=\sum_{l=P}^{H-1}a_{h-l}v_l.
\end{equation}
The set of all possible witnesses, even as~$s$ varies, is finite for each fixed~$H$: $\alpha_s$ lies between $\lambda$ and $\lambda/\rho$, and every prime and integer in \eqref{eq:tail-bands} is consequently bounded in terms of~$H$. Also $D(h,\eta)=O_H(\log(h+1))$. Thus the series used below converges absolutely, uniformly over this finite set and over~$s$.

A tail totient must contribute once even when it has several witnesses.
The union that enforces this can be described by independent exponential
random variables $E_h$ of mean one, indexed by $h\ge H$. For each witness
put
\begin{equation}\label{eq:tail-events}
 \mathcal E_\eta=
 \bigcap_{h\ge H}
 \left\{E_h\ge\frac\gamma{\alpha_s}\rho^hD(h,\eta)\right\}.
\end{equation}
For a function $f:[1,\infty)\to[0,1]$, the coefficient is
\[
 A_H(f;s)=\rho^{H(H-1)/2}
       \left(\frac\gamma{\alpha_s}\right)^H
       \sum_{\substack{d\in\mathcal V\\\mathcal W_{H,s}(d)\ne\varnothing}}
       \frac{f(\ell(d)/d)}d\,
       \Pr\left(\bigcup_{\eta\in\mathcal W_{H,s}(d)}
                    \mathcal E_\eta\right).
\]
An intersection of events requires each $E_h$ to exceed the largest
of the corresponding thresholds. Independence and absolute convergence
of their sum therefore turn finite inclusion--exclusion into the
explicit expression
\begin{equation}\label{eq:AH}
\begin{split}
 A_H(f;s)={}&\rho^{H(H-1)/2}
       \left(\frac{\gamma}{\alpha_s}\right)^H
       \sum_{\substack{d\in\mathcal V\\\mathcal W_{H,s}(d)\ne\varnothing}}
       \frac{f(\ell(d)/d)}d\\[-2pt]
 &\quad\times
 \sum_{\varnothing\ne T\subseteq\mathcal W_{H,s}(d)}
       (-1)^{|T|-1}
 \exp\left\{-\frac{\gamma}{\alpha_s}
       \sum_{h=H}^{\infty}\rho^h
                    \max_{\eta\in T}D(h,\eta)\right\}.
\end{split}
\end{equation}
For $f=1$, the factor involving~$\ell$ is simply~$1$: no least-preimage
search is part of the main coefficient. For other weights, any witness
with totient~$d$ supplies the finite search bound $\ell(d)\le w(\eta)$.
The probability expression also shows directly that $A_H(f;s)\ge0$.
Section~\ref{sec:limits} identifies these probabilities with limiting
normalized prefix volumes: $1/d$ comes from counting the largest prime, and the
prefactor is the limiting ratio of the prefix volume to~$G_m$.

When the limit exists, write
\[
 A(f;s)=\lim_{H\to\infty}A_H(f;s).
\]
Each approximant uses finite arithmetic data and an absolutely convergent
series, not an unspecified counting constant. Theorem~\ref{thm:main}
justifies the limit for $f=1$; the next result does so for the weights
needed to count least preimages. Convergence for arbitrary $f$ is not
required.

\subsection{Least-preimage intervals}
For a positive integer~$k$, define
\begin{equation}\label{eq:fk}
 \begin{aligned}
 N_k(x)&=\#\{v\in\mathcal V:v\le x,\ kx<\ell(v)\le(k+1)x\},\\
 f_k(r)&=\min\{1,(k+1)/r\}-\min\{1,k/r\}.
 \end{aligned}
\end{equation}
The function $f_k$ is nonnegative and bounded by~$1$.

\begin{theorem}\label{thm:companion}
For each fixed positive integer $k$, the functions $A_H(f_k;s)$ converge uniformly for $0\le s<1$, and
\[
 N_k(x)=\frac{x}{\log x}\,G_m\bigl(A(f_k;\theta)+o(1)\bigr).
\]
There are two possibilities.
\begin{enumerate}[label=\textup{(\roman*)}]
 \item If some $d\in\mathcal V$ satisfies $\ell(d)>kd$, then
 \[
  \inf_{0\le s<1}A(f_k;s)>0,\qquad
  N_k(x)\sim\frac{x}{\log x}\,G_m A(f_k;\theta)
       \asymp_k V(x).
 \]
 \item If no such $d$ exists, then $N_k(x)=0$ for every $x>0$, and $A(f_k;s)=0$ for every $s\in[0,1)$.
\end{enumerate}
The first alternative holds for $k=1,2$.
\end{theorem}

This theorem does not determine which other integers~$k$ satisfy the first alternative. Whether the first alternative holds for every positive integer~$k$ is equivalent to the unboundedness of $\ell(d)/d$; the weighted formula does not settle this question.

\subsection{Proof strategy}
Two parts of the argument ensure that this construction counts values once. First, outside an exceptional set, \emph{every} preimage has the required form. Second, collisions between different long prefixes have negligible total multiplicity. For the latter we use the layered shifted-prime method of
Maier and Pomerance \cite[Section~3]{MaierPomerance1988}, in the form
developed by Ford, and adapt his comparison \cite[Lemma~5.1]{FordTotients2013} to a different set of hypotheses, proving the necessary uniform bound in Section~\ref{sec:layers}. The number of layers grows, but each sieve application concerns only two or three linear forms. The comparison and the subsequent control of tuple multiplicities are the ingredients that turn a volume estimate into an asymptotic for distinct values.

The end of the continuous prefix and the end of the retained prime list
are different cuts. Their separation makes the exceptional-set estimates
small enough to absorb the smooth tail. Section~\ref{sec:extraction}
obtains the required preimages directly from Ford's Theorems~10 and~16
and develops the volume estimates. Section~\ref{sec:collisions} proves
prefix uniqueness and transfers tuple counts to distinct values.
Section~\ref{sec:limits} counts the largest prime using common endpoint
parameters, obtaining fixed dilation before identifying the arithmetic
limit. Finally, Section~\ref{sec:companion} weights that prime count to
locate least preimages and proves the positive and zero alternatives.
